\documentclass[11pt]{article}
\newcommand{\courseshort}{Data Structures (Make-up) · 010153103}
\newcommand{\lecturetitle}{L6: Trees \& Binary Trees}
\usepackage{lecture_style}

\begin{document}
\lectureheader{Lecture 6: Trees \& Binary Trees}{Pre-class brief}{Goodrich Ch.~8}

\begin{objbox}
\begin{itemize}[leftmargin=*]
  \item Define tree terminology: root, child, parent, leaf, depth, height, ancestor.
  \item Implement a linked binary-tree node with \texttt{left}, \texttt{right}, \texttt{parent}.
  \item Perform preorder, inorder, postorder, and level-order traversals.
  \item Compute the height of a tree recursively.
\end{itemize}
\end{objbox}

\section*{1. Vocabulary}

\begin{itemize}[leftmargin=*]
  \item \textbf{Root}: the topmost node (no parent).
  \item \textbf{Leaf}: a node with no children.
  \item \textbf{Depth} of a node: number of edges from the root to that node.
  \item \textbf{Height} of a tree: the maximum depth of any node (a single-node tree has height 0).
  \item \textbf{Subtree} rooted at $v$: $v$ and all its descendants.
\end{itemize}

\section*{2. Binary Tree --- linked node}

\begin{lstlisting}
class TNode<E> {
    E element;
    TNode<E> left, right, parent;
    TNode(E e) { element = e; }
}
\end{lstlisting}

A binary tree is \textbf{full} if every node has 0 or 2 children, \textbf{complete} if all levels are full except possibly the last, which is filled left-to-right, and \textbf{perfect} if it is full and all leaves have the same depth.

\section*{3. The Four Traversals}

\begin{center}
\begin{tikzpicture}[
  every node/.style={circle, draw, minimum size=0.7cm, font=\small\sffamily}
]
  \node (a) {A};
  \node[below left=0.8cm and 1.3cm of a] (b) {B};
  \node[below right=0.8cm and 1.3cm of a] (c) {C};
  \node[below left=0.8cm and 0.4cm of b] (d) {D};
  \node[below right=0.8cm and 0.4cm of b] (e) {E};
  \node[below right=0.8cm and 0.4cm of c] (f) {F};
  \draw (a) -- (b); \draw (a) -- (c);
  \draw (b) -- (d); \draw (b) -- (e); \draw (c) -- (f);
\end{tikzpicture}
\end{center}

For the tree above:
\begin{itemize}[leftmargin=*]
  \item \textbf{Preorder} (root, left, right): A, B, D, E, C, F
  \item \textbf{Inorder} (left, root, right): D, B, E, A, C, F
  \item \textbf{Postorder} (left, right, root): D, E, B, F, C, A
  \item \textbf{Level-order} (BFS, uses a queue): A, B, C, D, E, F
\end{itemize}

\begin{lstlisting}
void preorder(TNode<E> v) {
    if (v == null) return;
    visit(v);
    preorder(v.left);
    preorder(v.right);
}
\end{lstlisting}

\section*{4. Useful recursive computations}

\begin{lstlisting}
int height(TNode<E> v) {
    if (v == null) return -1;
    return 1 + Math.max(height(v.left), height(v.right));
}
int size(TNode<E> v) {
    if (v == null) return 0;
    return 1 + size(v.left) + size(v.right);
}
\end{lstlisting}

\begin{exbox}{Pre-Class Exercises (do BEFORE the lecture)}
\textbf{Exercise 6.1 --- Build by hand.} Draw the binary tree obtained by inserting (in order) the keys
\verb|50, 30, 70, 20, 40, 60, 80| as a binary tree where smaller goes left and larger goes right (a BST). Then list its pre-, in-, post-, and level-order traversals.

\textbf{Exercise 6.2 --- Implement.} Write recursive Java methods (with the signature you need):
\begin{enumerate}[label=(\alph*)]
  \item \verb|countLeaves(v)| --- number of leaf nodes in the subtree rooted at \texttt{v}.
  \item \verb|sumElements(v)| --- sum of all integer elements (assume \verb|TNode<Integer>|).
  \item \verb|depth(v)| --- the depth of node \texttt{v} (use the parent pointer; iterative is fine).
\end{enumerate}

\textbf{Exercise 6.3 --- Reconstruct.} Given preorder $=$ A,B,D,E,C,F and inorder $=$ D,B,E,A,C,F, reconstruct the tree. (You'll discuss the algorithm in class.)

\textbf{Exercise 6.4 --- Level order.} Implement \verb|levelOrder(root)| using a \verb|Queue<TNode<E>>|. Print the elements in level order.

\textbf{Exercise 6.5 --- True or False (justify).}
\begin{enumerate}[label=(\alph*)]
  \item In a perfect binary tree of height $h$, the number of leaves is $2^h$.
  \item In a binary tree with $n$ nodes, height can be as small as $\lfloor\log_2 n\rfloor$ and as large as $n-1$.
  \item Inorder traversal of a BST visits keys in sorted order.
\end{enumerate}

\textbf{Exercise 6.6 --- Mirror.} Write a method that mirrors a binary tree (swap left/right of every node) in $O(n)$ time.
\end{exbox}

\begin{disbox}
\begin{itemize}[leftmargin=*]
  \item Why is preorder ``natural'' for printing structure but inorder ``natural'' for sorted output?
  \item What's the connection between traversal order and the data structures used (stack vs.\ queue)?
  \item What goes wrong with \texttt{height} if the tree is very deep? (Preview: stack vs.\ heap memory.)
\end{itemize}
\end{disbox}

\end{document}
